%==============================================================================
% Sjabloon poster bachproef
%==============================================================================
% Gebaseerd op document class `a0poster' door Gerlinde Kettl en Matthias Weiser
% Aangepast voor gebruik aan HOGENT door Jens Buysse en Bert Van Vreckem

\documentclass[a0,portrait]{hogent-poster}

% Info over de opleiding
\course{Bachelorproef}
\studyprogramme{toegepaste informatica}
\academicyear{2025-2026}
\institution{Hogeschool Gent, Valentin Vaerwyckweg 1, 9000 Gent}

% Info over de bachelorproef
\title{De populariteit van point-of-interests voorspellen op basis van censusdata met neurale netwerken.}
\author{Ruben Priels}
\email{ruben.priels@student.hogent.be}
\supervisor{Dhr. K. Mertens}
\cosupervisor{Dhr. M. Roels (Accurat)}

% Indien ingevuld, wordt deze informatie toegevoegd aan het einde van de
% abstract. Zet in commentaar als je dit niet wilt.
\specialisation{AI \& data engineer}
\keywords{POI, Neurale netwerken, MLP, RNN, Transformer}
% \projectrepo{https://github.com/user/repo}

\begin{document}

\maketitle

% \begin{abstract}
% Veel bedrijven beschikken vaak over onvoldoende inzicht in hoe % aantrekkelijk een locatie is voor klanten, wat leidt tot slechte % vestigingskeuzes en foutieve inkoopbeslissingen. Hierbij werd dus % onderzocht of artificiële intelligentie de populariteit van een % POI kan voorspellen op basis van censusdata. In de proof of % concept werden drie neurale netwerken (MLP, RNN en transformer) % getraind en geëvalueerd. De transformer presteerde het beste op % zowel reguliere dagen als feestdagen en leverde bruikbare % voorspellingen op tot minstens een maand in de toekomst. De studie % concludeert dat neurale netwerken een waardevolle meerwaarde % bieden voor bedrijven bij strategische locatie- en % inkoopbeslissingen.
% \end{abstract}

\begin{multicols}{2} % This is how many columns your poster will be broken into, a portrait poster is generally split into 2 columns

\section{Probleemstelling}

Bedrijven die afhankelijk zijn van fysieke locaties beschikken vaak over onvoldoende inzicht in de aantrekkelijkheid van een locatie voor zowel huidige als potentiële klanten. Dit leidt tot verkeerde beslissingen bij het openen van nieuwe vestigingen of het inkopen van voorraden. Wanneer een onderneming zich vestigt op een locatie met lage bezoekersaantallen, leidt dit tot een lager klantenaantal en bijgevolg een lagere omzet. Daarnaast kan een onvoldoende inzicht in de populariteit van een POI leiden tot foutieve inkoopbeslissingen, met als gevolg hoge voorraadkosten en verspilling of gemiste verkopen. Door een beter inzicht te krijgen in de toekomstige populariteit van locaties kunnen ondernemingen niet alleen strategische beslissingen nemen bij het openen van nieuwe vestigingen, maar ook nauwkeuriger inschatten hoeveel goederen aangekocht moeten worden voor bestaande POI’s.

\section{Onderzoeksvraag}

Om dit probleem aan te pakken, richtte deze studie zich op de onderzoeksvraag: “Hoe kan artificiële intelligentie worden ingezet om de populariteit van een point-of-interest te voorspellen op basis van census data?” Deze onderzoeksvraag werd opgesplitst in vier deelvragen.

\vspace{5mm}

\begin{enumerate}
    \item Wat betekent de populariteit van een POI en hoe kan deze gemeten worden?
    \item Hoe wordt de populariteit van een POI verklaard door censusdata?
    \item Welke neurale netwerken zijn geschikt om POI-populariteit te voorspellen?
    \item Welke technieken bestaan er om de prestaties van neurale netwerken te verbeteren?
\end{enumerate}

\section{Methodologie}

Het onderzoek werd verdeeld in vijf verschillende fasen.

\textbf{De literatuurstudie:} Onderzoek naar bestaande literatuur om de hoofd- en deelvragen op te lossen.
    
\textbf{Dataverzameling:} Het verzamelen en samenvoegen van diverse databronnen tot een betrouwbare dataset.
    
\textbf{Implementatie neurale netwerken:} Het ontwikkelen en trainen van drie verschillende types neurale netwerken.
    
\textbf{Experimenten uitvoeren:} Diverse experimenten uitvoeren om de prestaties van de modellen te meten.
    
\textbf{Evaluatie van de neurale netwerken:} Het analyseren en vergelijken van de resultaten om te bepalen welk model het meest geschikt is.

\vspace{10mm}

\noindent
\begin{minipage}{\linewidth}
    \centering
    
    \includegraphics[width=0.4\linewidth]{methodologie_diagram}
    
    \captionof{figure}{Fasen van het onderzoek}
\end{minipage}

\section{Experimenten}
\setlength{\tabcolsep}{15pt}
\renewcommand{\arraystretch}{1.2}
%\textbf{Globale prestaties (alle dagen)}\\[2mm]

\vspace{10mm}

De algemene prestaties van de verschillende modellen werden vergeleken om te bepalen welk model over het algemeen het best presteert. Daarnaast werd de meerwaarde van neurale netwerken aangetoond.

\begin{center}
\begin{tabular}{|l|ccc|}
\toprule
\textbf{Model} & \textbf{MAE} & \textbf{RMSE} & \textbf{R²} \\
\midrule
            Naïeve & 148.5 & 195.2 & 0.32 \\
            XGBoost             & 124.0 & 169.8 & 0.45 \\
            MLP                 & 123.96 & 169.46 & 0.47 \\
            RNN                 & 112.77 & 153.34 & 0.52 \\
            \textbf{Transformer}         & \textbf{105.12} & \textbf{144.56} & \textbf{0.55} \\
            \bottomrule
        \end{tabular}
        \captionof{table}{Globale prestaties (alle dagen)}
\end{center}

\vspace{10mm}

% hier miss enkel transformer?
De prestaties van elk model op reguliere dagen en feestdagen werden vergeleken. Op deze manier werd geëvalueerd hoe nauwkeurig de modellen voorspellen wanneer het bezoekersgedrag verandert.

% vervangen door grafiek van transformer van MAE en RMSE?
\begin{center}
\begin{tabular}{|l|ccc|ccc|}
            \toprule
            \textbf{Model} & \multicolumn{3}{c}{\textbf{Normaal}} &
            \multicolumn{3}{c}{\textbf{Feestdagen}} \\
            & MAE & RMSE & R\textsuperscript{2} & MAE & RMSE & R² \\
            \midrule
            MLP         & 109.03 & 146.23 & 0.58 & 155.49 & 214.18 & 0.36 \\
            RNN         & 99.12 & 134.96 & 0.62 & 145.77 & 199.84 & 0.41 \\
            \textbf{Transformer} & \textbf{91.37} & \textbf{126.42} & \textbf{0.66} & \textbf{138.84} & \textbf{191.53} & \textbf{0.45} \\
            \bottomrule
        \end{tabular}
        \captionof{table}{Reguliere dagen en feestdagen}
\end{center}

\vspace{10mm}

%\textbf{Voorspellingshorizon}\\[2mm]
De prestaties van de sequentiële modellen werden vergeleken voor verschillende voorspellingshorizonten. Op deze manier werd geëvalueerd hoe nauwkeurig de modellen voorspellingen maken op verschillende termijnen.

\begin{center}
        \begin{tabular}{|l|ccc|ccc|}
            \toprule
            \textbf{Termijn} & \multicolumn{3}{c}{\textbf{Transformer}} &
            \multicolumn{3}{c}{\textbf{RNN}} \\
            & MAE & RMSE & R\textsuperscript{2} & MAE & RMSE & R² \\
            \midrule
            Kort ($\leq$1 week)   & 103.08 & 141.29 & 0.57 & 111.2 & 149.72 & 0.53 \\
            Middel (1--3 weken)   & 106.02 & 145.08 & 0.56 & 113.29 & 153.86 & 0.51 \\
            Lang (3 weken--1 maand) & 106.47 & 147.03 & 0.55 & 113.82 & 154.89 & 0.51 \\
            \bottomrule
        \end{tabular}
        \captionof{table}{Voorspellingshorizon}
\end{center}

\vspace{10mm}

\section{Conclusies}

Deze studie toont aan dat neurale netwerken ingezet kunnen worden om de populariteit van een POI te voorspellen op basis van censusdata. De transformer presteert duidelijk het best. Het behaalt de laagste fouten en de hoogste verklaarde variantie, zowel op normale dagen als op feestdagen. Hoewel de voorspellingsfouten toenemen naarmate de tijd vordert, blijven deze afwijkingen beperkt waardoor de voorspellingen bruikbaar blijven tot minstens een maand in de toekomst.

Deze resultaten hebben invloed op verschillende aspecten van de bedrijfswereld. Bedrijven kunnen met dit model gerichter goederen inkopen en betere strategische beslissingen nemen bij het openen van een vestiging. Hierdoor zullen ondernemingen meer omzet kunnen maken en startende bedrijven zullen een grotere kans hebben om te overleven.

\section{Toekomstig onderzoek}

Dit onderzoek leidt tot nieuwe vragen. Hoe nauwkeurig kunnen de bezoekersaantallen voorspeld worden in extreme situaties zoals epidemieën of pandemieën zoals corona? Op welke manieren zou het model hier beter op kunnen voorspellen? Hoe kunnen de voorspellingen eenvoudiger verklaard worden voor bedrijven?

\end{multicols}
\end{document}